\documentclass[12pt,a4paper]{article}

\usepackage[spanish,es-tabla]{babel}
\usepackage[T1]{fontenc}
\usepackage[utf8]{inputenc}
\usepackage{lmodern}
\usepackage[a4paper,margin=2.5cm]{geometry}
\usepackage{setspace}
\usepackage{parskip}
\usepackage{hyperref}
\usepackage{graphicx}
\usepackage{float}
\usepackage{enumitem}

\hypersetup{
    colorlinks=true,
    linkcolor=blue,
    urlcolor=blue,
    pdftitle={Anexo Específico de la Instalación de Agua},
    pdfauthor={OpenAI Codex}
}

\setstretch{1.1}
\setlist[itemize]{topsep=4pt,itemsep=2pt}
\setlist[enumerate]{topsep=4pt,itemsep=2pt}

\newcommand{\placeholderimage}[2]{
    \begin{figure}[H]
        \centering
        \fbox{\parbox[c][6cm][c]{0.88\textwidth}{
            \centering
            \textbf{#1}\\[0.4cm]
            #2
        }}
        \caption{#2}
    \end{figure}
}

\newcommand{\manualimage}[2]{
    \IfFileExists{#1}{
        \begin{figure}[H]
            \centering
            \includegraphics[width=0.9\textwidth]{#1}
            \caption{#2}
        \end{figure}
    }{
        \placeholderimage{#1}{#2}
    }
}

\title{Anexo Específico de la Instalación de Agua}
\author{}
\date{\today}

\begin{document}

\maketitle
\tableofcontents
\newpage

\section{Finalidad de este anexo}

Este anexo complementa el manual global y recoge únicamente las particularidades de la instalación de Agua.

El objetivo es que el usuario sepa qué elementos, sistemas y reglas concretas debe tener en cuenta cuando trabaja con Agua, sin mezclar esta información con el funcionamiento común del resto de instalaciones.

\section{Qué se documenta aquí}

En este anexo se explican:

\begin{itemize}
\item los sistemas disponibles dentro de Agua;
\item las distribuciones utilizadas;
\item los diámetros habituales;
\item los accesorios automáticos;
\item los elementos definidos propios de Agua;
\item y las limitaciones o avisos que el usuario debe conocer.
\end{itemize}

\section{Sistemas disponibles en Agua}

En la instalación de Agua, el usuario puede trabajar con estos sistemas:

\begin{itemize}
\item Polietilè;
\item Multicapa;
\item Armaflex.
\end{itemize}

Cada uno de ellos pertenece a la misma instalación general, pero representa una solución distinta dentro del modelo.

\manualimage{CAPTURA-AGUA-01.png}{Desplegable con los sistemas disponibles en Agua}

\section{Distribuciones en Agua}

En Agua el usuario trabaja principalmente con estas distribuciones:

\begin{itemize}
\item TD;
\item IS.
\end{itemize}

La distribución condiciona cómo se comporta la instalación y qué configuración se aplica al resultado final.

\subsection{Qué debe tener en cuenta el usuario}

\begin{itemize}
\item La distribución debe elegirse antes de dibujar.
\item Si se cambia después, conviene revisar el resultado.
\item No todos los sistemas se comportan igual en todas las distribuciones.
\end{itemize}

\subsection{Caso especial}

En la configuración actual, si el usuario intenta trabajar con combinaciones no admitidas en ciertos sistemas, la herramienta muestra un aviso para corregir la configuración.

Conviene revisar especialmente los casos de Multicapa y Armaflex cuando la distribución elegida no sea la adecuada.

\manualimage{CAPTURA-AGUA-02.png}{Campo de distribución y ejemplo de aviso por combinación no admitida}

\section{Diámetros en Agua}

En la instalación de Agua, los diámetros más habituales en la configuración actual son:

\begin{itemize}
\item 20;
\item 25.
\end{itemize}

El diámetro afecta a:

\begin{itemize}
\item el tamaño del tubo;
\item los accesorios compatibles;
\item la geometría final;
\item y parte de la información asociada al elemento.
\end{itemize}

\subsection{Recomendación}

Si se modifica el diámetro después de dibujar, conviene revisar:

\begin{itemize}
\item uniones;
\item codos;
\item bifurcaciones;
\item y el resultado final en pantalla.
\end{itemize}

\section{Accesorios automáticos en Agua}

Agua genera automáticamente distintos accesorios en función del recorrido.

\subsection{Codos}

En Agua, el usuario debe tener presente que el sistema trabaja principalmente con codos de 90 grados.

Esto significa que el recorrido debe plantearse teniendo en cuenta esa lógica de giro.

\subsection{Uniones o manguitos}

Se generan cuando dos tramos deben unirse en línea.

\subsection{Bifurcaciones}

Se generan cuando el recorrido lo requiere y la geometría del trazado lo justifica.

\manualimage{CAPTURA-AGUA-03.png}{Ejemplo real de un codo de 90 grados generado en Agua}

\section{Qué diferencia a Agua de otras instalaciones}

Una diferencia importante respecto a otras instalaciones es que Agua debe explicarse con sus propias reglas de uso.

Por ejemplo:

\begin{itemize}
\item en Agua hay sistemas como Polietilè, Multicapa y Armaflex;
\item en otras instalaciones puede haber familias completamente distintas;
\item en Agua el usuario debe pensar en codos de 90 grados;
\item mientras que en otras instalaciones, como Saneamiento, puede haber además codos de 45 grados.
\end{itemize}

Por eso no conviene meter toda esta información específica dentro del manual global. Es mejor explicarla en anexos separados.

\section{Elementos definidos propios de Agua}

Además del recorrido principal, Agua permite insertar elementos definidos propios del sistema.

En la configuración actual, Agua trabaja con elementos definidos como:

\begin{itemize}
\item Colze Base;
\item Clau de Pas;
\item Taps;
\item T sortida.
\end{itemize}

\subsection{Qué debe entender el usuario}

Estos elementos no forman parte del trazado automático básico, sino que se colocan cuando el usuario necesita completar la instalación con piezas específicas.

\manualimage{CAPTURA-AGUA-04.png}{Paleta de Elementos Definidos en Agua}

\section{Macros en Agua}

Agua también permite insertar macros como complemento del recorrido.

Se recomienda:

\begin{itemize}
\item dibujar primero el recorrido principal;
\item revisar la geometría;
\item y después insertar las macros necesarias.
\end{itemize}

\section{Layers y atributos en Agua}

Agua utiliza sus propios layers y parte de su propia lógica de atributos.

\subsection{Qué debe saber el usuario}

\begin{itemize}
\item No todos los sistemas aplican exactamente la misma información.
\item La distribución influye en el comportamiento del resultado.
\item Algunos elementos pueden conservar configuraciones específicas del propio modelo.
\end{itemize}

Lo importante es comprobar siempre:

\begin{itemize}
\item que el layer aplicado es el correcto;
\item que la distribución es la adecuada;
\item y que la instalación final se corresponde con el sistema seleccionado.
\end{itemize}

\section{Puntos a revisar con especial atención}

En Agua conviene revisar especialmente estos puntos:

\subsection{Sistema seleccionado}

Comprobar siempre si se está trabajando en Polietilè, Multicapa o Armaflex.

\subsection{Distribución}

Comprobar si el trabajo corresponde realmente a TD o IS.

\subsection{Giros del recorrido}

Tener presente que la lógica habitual de Agua se apoya en giros de 90 grados.

\subsection{Resultado después de un cambio}

Si se cambia diámetro, sistema o distribución, conviene revisar el recorrido antes de finalizar.

\section{Cómo replicar esta documentación en otras instalaciones}

La mejor forma de documentar el conjunto completo no es ampliar indefinidamente el manual global, sino crear un anexo específico para cada instalación.

La estructura recomendada para cada instalación es esta:

\begin{enumerate}
\item Sistemas disponibles.
\item Distribuciones disponibles.
\item Diámetros.
\item Accesorios automáticos.
\item Elementos definidos propios.
\item Macros, si aplica.
\item Layers y atributos a vigilar.
\item Limitaciones o avisos especiales.
\item Casos prácticos o ejemplos visuales.
\end{enumerate}

\section{Resumen final}

El manual global ya explica correctamente el funcionamiento común.

Lo que faltaba era separar las particularidades de cada instalación en anexos independientes. Este documento resuelve esa necesidad para Agua y puede usarse como modelo para crear otros anexos, por ejemplo:

\begin{itemize}
\item anexo de Saneamiento;
\item anexo de Ventilación;
\item anexo de Electricidad.
\end{itemize}

\end{document}
